%! Operations with Polynomials — Unit 4, Lesson 4.1 — Homework
\documentclass[10pt]{article}
\usepackage{algebra2-article}
\usepackage{algebra2-boxes}
\newcommand{\TallMath}[1]{$\displaystyle #1\rule[-1.4em]{0pt}{3.2em}$}

\begin{document}

\pageheader{Unit 4, Lesson 4.1}{Homework}
\namedateperiod

\vspace{0.10in}

\begin{notesbox}{Practice}
Write every answer in \textbf{standard form}. Show the products or the sign changes you used.

\begin{enumerate}[leftmargin=1.9em, itemsep=11pt]
  \item \textbf{Name it.} Consider $8 - 3x^2 + x^5 - 6x$.
    \begin{enumerate}[label=(\alph*), leftmargin=1.8em, itemsep=4pt]
      \item Standard form: \blank{6.0cm}
      \item Degree: \blank{1.0cm} \quad Leading coefficient: \blank{1.0cm} \quad Number of terms:
            \blank{1.0cm}
      \item Based on the degree and leading coefficient, both arms of its graph: \blank{3.6cm}
    \end{enumerate}

  \item \textbf{Add and subtract.}
    \begin{enumerate}[label=(\alph*), leftmargin=1.8em, itemsep=6pt]
      \item $(6x^3 - x + 4) + (2x^2 + 5x - 9) = $ \blank{5.5cm}
      \item $(9x^2 - 3x + 2) - (4x^2 - 3x - 6) = $ \blank{5.5cm}
    \end{enumerate}

  \item \textbf{Multiply.}
    \begin{enumerate}[label=(\alph*), leftmargin=1.8em, itemsep=6pt]
      \item $-3x^2(2x^3 - x + 5) = $ \blank{5.5cm}
      \item $(4x - 1)(x + 7) = $ \blank{5.5cm}
      \item $(x - 2)(x^2 + 3x - 5) = $ \blank{5.5cm}
    \end{enumerate}

  \item \textbf{Special products.} Use the identities --- no box needed.

    \vspace{3pt}
    (a) $(x + 9)^2 = $ \blank{4.3cm} \qquad (b) $(2x - 5)^2 = $ \blank{4.3cm}

    \vspace{4pt}
    (c) $(3x + 8)(3x - 8) = $ \blank{3.8cm} \qquad (d) $(x + 6y)(x - 6y) = $ \blank{3.8cm}

  \item \textbf{Back to a graph you already read.} On yesterday's Exit Ticket you read the graph of
        $k(x) = (x+1)^2(x-2)$.
    \begin{enumerate}[label=(\alph*), leftmargin=1.8em, itemsep=5pt]
      \item Expand $(x+1)^2$: \blank{4.2cm}
      \item Now multiply by $(x-2)$ and write $k$ in standard form: \blank{5.5cm}
      \item Degree: \blank{1.0cm} \quad Leading coefficient: \blank{1.0cm} \quad End behavior:
            \blank{3.4cm}
      \item Substituting $x=0$ into your standard form gives the $y$-intercept \blank{1.8cm}. Does it
            agree with the graph you read yesterday? \blank{1.4cm}
    \end{enumerate}

  \item \textbf{Explain.} Why does subtracting a polynomial require changing the sign of \emph{every}
        term inside the parentheses, and not just the first one? Support your answer with a short
        example that goes wrong if you change only the first sign. \writelines{3}
\end{enumerate}
\end{notesbox}

\begin{extensionbox}
\textbf{An area model that needs both operations.} A rectangular garden measures $(2x+3)$ feet by
$(x+5)$ feet. A square pond with side $x$ feet is dug inside it, and the rest is planted.
\begin{enumerate}[label=(\alph*), leftmargin=1.8em, itemsep=5pt]
  \item Write and expand a polynomial for the \emph{total} garden area.
  \item Write a polynomial for the \emph{planted} area (total minus pond) in standard form.
  \item Use your answer to (b) to find the planted area when $x = 4$ feet.
  \item Check part (c) a second way: find the garden's actual dimensions and the pond's actual area at
        $x=4$, and subtract. Do the two methods agree?
\end{enumerate}
\writelines{4}
\end{extensionbox}

\begin{spiralbox}
\textbf{Coming next --- Lesson 4.2: Advanced Factoring.} Today you ran products \emph{forward}:
$(x-1)^2(x+2)$ became $x^3-3x+2$. Next you will run them \emph{backward} --- pulling out a GCF,
factoring by grouping, and recognizing the difference of squares and the new sum and difference of
\emph{cubes} --- turning a standard form back into the factored form that shows you the zeros.
\end{spiralbox}

\end{document}
